
The coefficient of variation of participation ratio (CV\_PR), hypothesized to correlate negatively with ImageNet accuracy, instead shows a strong positive correlation (r = +0.6065, p = 9.24 $\times 10^{-11})$. Models with higher accuracy exhibit more variance in randomized SVD estimates, not less. The 95\% confidence interval [0.506, 0.703] excludes all negative values.

\subsubsection{Summary of Contributions}

\begin{enumerate}
\item \textbf{Methodology validation:} CV\_PR can be reliably extracted from 100 pretrained models using randomized SVD with 20 seeds (100\% completion, CV\_PR in [0.0014, 0.0326]).
\end{enumerate}

\begin{enumerate}
\item \textbf{Hypothesis falsification:} The negative correlation hypothesis is rejected. CV\_PR positively correlates with accuracy.
\end{enumerate}

\begin{enumerate}
\item \textbf{Interpretive framework:} Confounding by model size and richer representations are identified as competing explanations.
\end{enumerate}

\subsubsection{Future Directions}

\textbf{Confound analysis:} Partial correlation controlling for parameter count would distinguish whether CV\_PR has independent predictive value.

\textbf{Mechanism investigation:} The question of why higher spectral variance accompanies better generalization warrants investigation with reformulated hypotheses.

\textbf{Architecture stratification:} Within-family versus cross-family analysis may reveal whether the positive correlation varies by architecture type.
